\newpage

\appendix
\section{Prompts Used in Feature Extraction}
\label{appendix:prompts}

This appendix details the system prompts provided to the \textit{Gemma-2-9b-it} model. The structure of these prompts follows the model's control tag format to ensure a deterministic response in JSON format.

\subsection{Subtopic Extraction}
\label{appendix:prompt_subtopics}

Experiment C constitutes the granular foundation of the system. Two variants were used: one based on a rubric (0-5) and a binary version (0/1).

\begin{lstlisting}[language=json, caption={Prompt for Subtopics rubric  (Experiment C).}, label={lst:app_prompt_sub_rubric}, frame=single, basicstyle=\small\ttfamily, breaklines=true]
<start_of_turn>user
You are a text analysis system. Your goal is to quantify the level of self-disclosure in a predefined list of topics:
- Health (Mental/Physical)
- Relationships (Family/Affective)
- Personal History (Goals, Traumas, Career)

Analyze the [MESSAGES TO EVALUATE] and assign a score (0-5) to EACH topic. 
Response must be ONLY a valid JSON object.

# EVALUATION SCALE (0-5)
* 0: No Mention.
* 1: Impersonal Mention (General talk).
* 2: Vague Personal Mention (General identification).
* 3: Specific Action Disclosure (Concrete habits/actions).
* 4: Detailed Disclosure (Sharing reasoning/beliefs).
* 5: Deep Disclosure (Vulnerability, impact, deep emotions).

# EXPECTED JSON FORMAT
{
  "mental_emotional_health": Int,
  "physical_health_addictions": Int,
  "family_dynamics": Int,
  "affective_relationships": Int,
  "personal_goals_failures": Int,
  "traumas_adversities": Int,
  "professional_career": Int
}

[MESSAGES TO EVALUATE]: 
<end_of_turn>
\end{lstlisting}

\begin{lstlisting}[language=json, caption={Prompt for Subtopics binary (Experiment C).}, label={lst:app_prompt_sub_binary}, frame=single, basicstyle=\small\ttfamily, breaklines=true]
<start_of_turn>user
You are a high-precision text classification system. Your task is to detect the PRESENCE or ABSENCE of specific self-disclosure topics in the provided text.

Analyze the [MESSAGES TO EVALUATE] and assign a value of 1 if the topic is mentioned (regardless of its depth) and 0 if it is not mentioned at all.
Your response must be ONLY a valid JSON object.

# TOPICS TO DETECT
- mental_emotional_health
- physical_health_addictions
- family_dynamics
- affective_relationships
- personal_goals_failures
- traumas_adversities
- professional_career

# DETECTION RULES
* 0: There is no trace of the topic in the text.
* 1: The topic is present in the text (even if vaguely mentioned).

# EXPECTED JSON FORMAT
{
  "mental_emotional_health": 0 or 1,
  "physical_health_addictions": 0 or 1,
  "family_dynamics": 0 or 1,
  "affective_relationships": 0 or 1,
  "personal_goals_failures": 0 or 1,
  "traumas_adversities": 0 or 1,
  "professional_career": 0 or 1
}

[MESSAGES TO EVALUATE]:
<end_of_turn>
\end{lstlisting}

\newpage

\subsection{Direct topic Evaluation}
\label{appendix:prompt_categories}

In Experiment B, the model performs a direct inference on the three macro-categories of the SPT, bypassing the aggregation of subtopics.

\begin{lstlisting}[language=json, caption={Prompt for Topic rubric (Experiment B).}, label={lst:app_prompt_cat_rubric}, frame=single, basicstyle=\small\ttfamily, breaklines=true]
<start_of_turn>user
You are a text analysis system. Your goal is to quantify the level of self-disclosure in a predefined list of categories focused on lived experiences and the author's inner world.

Analyze the [MESSAGES TO EVALUATE] and assign a score (0-5) to EACH category.
Response must be ONLY a valid JSON object.

# CATEGORIES TO EVALUATE
* 'Category 1: Health and Well-being': Includes physical, mental, and emotional health.
* 'Category 2: Interpersonal Relationships': Includes family dynamics and affective relationships.
* 'Category 3: Personal History and Development': Includes goals, failures, traumas, and career.

# EVALUATION SCALE (0-5)
* 0: No Mention.
* 1: Impersonal Mention (General discussion, no personal link).
* 2: Vague Personal Mention (General identification or labels).
* 3: Specific Action Disclosure (Concrete habits or specific actions)
* 4: Detailed Disclosure (Reasoning, beliefs, or specific data).
* 5: Deep Disclosure (Vulnerability, fear, or profound personal impact).

# EXPECTED JSON FORMAT
{
    "Category 1: Health and Well-being": Int,
    "Category 2: Interpersonal Relationships": Int,
    "Category 3: Personal History and Development": Int
}

[MESSAGES TO EVALUATE]: 
<end_of_turn>
\end{lstlisting}

\begin{lstlisting}[language=json, caption={Prompt for topic binary (Experiment B).}, label={lst:app_prompt_cat_binary}, frame=single, basicstyle=\small\ttfamily, breaklines=true]
<start_of_turn>user
You are a text analysis system. Your goal is to quantify the level of self-disclosure in a predefined list of broad categories:
- Category 1: Health and Well-being (Mental, emotional, and physical health)
- Category 2: Interpersonal Relationships (Family dynamics, affective relationships).
- Category 3: Personal History and Development (Goals, failures, traumas, career).

Analyze the [MESSAGES TO EVALUATE] and assign a 0 or 1 to EACH category.
Response must be ONLY a valid JSON object.

# BINARY EVALUATION SCALE (0-1)
* 0: No Mention or Not Relevant: The topic is not mentioned or not linked personally.
* 1: Specific Action/Deep Disclosure: The author reveals a concrete habit, action, or reflects deeply on the topic.

# EXPECTED JSON FORMAT
{
    "Category 1: Health and Well-being": 0 or 1,
    "Category 2: Interpersonal Relationships": 0 or 1,
    "Category 3: Personal History and Development": 0 or 1
}

[MESSAGES TO EVALUATE]: 
<end_of_turn>
\end{lstlisting}

\newpage
